When dealing with systems of PDEs (as opposed to single equations) we will not, in general, focus on a more restrictive definition of ellipticity. That is, we will stick to the definition of an elliptic system as one with no real characteristics, and characteristics are to be determined by the principal part of the operator defined using appropriate weights. However, the reader should be aware of some particularly important examples of elliptic systems that arise in the calculus of variations. Let $\Omega$ be a domain in $\mathbb{R}^n$ and let $\bm{u}:\Omega\to\mathbb{R}^N$. Consider the system of $N$ second-order differential equations in divergence form
\begin{equation}\label{eq:second-order-divergence-system-definition}
\frac{\partial}{\partial x_k}\left(A_{kl}^{IJ}(\bm{x})\frac{\partial u_J}{\partial x_l}\right)+c^{IJ}(\bm{x})u_J(\bm{x})=0,\quad I=1,\ldots,N.
\end{equation}


Here, the coefficients $A_{kl}^{IJ}$ and $c^{IJ}$, $I,J=1,\ldots,N$, $k,l=1,\ldots,n$, are assumed to be sufficiently smooth, and the summation convention is assumed to hold from $1$ to $N$ for repeated uppercase indices and from $1$ to $n$ for repeated lowercase indices.

Note that if
\begin{equation}\label{eq:elliptic-system-full-matrix-positivity}
M_k^IA_{kl}^{IJ}M_l^J>0
\end{equation}
for every $\bm{x}\in\Omega$ and for every nonzero $N\times n$ matrix $\bm{M}$, then we can show that the system is elliptic simply by taking the ``obvious\textquotesingle\textquotesingle\ principal part (giving weight $1$ to each of the equations and dependent variables). Condition \eqref{eq:elliptic-system-full-matrix-positivity} or the uniform version
\begin{equation}\label{eq:elliptic-system-full-matrix-uniform-bound}
M_k^IA_{kl}^{IJ}M_l^J\geq\theta|\bm{M}|^2
\end{equation}
for every $\bm{x}\in\Omega$ and every nonzero $N\times n$ matrix $\bm{M}$ for some $\theta>0$ is often given as a definition of an elliptic system. However, such a definition does not fit such systems as the Stokes system.

Another important ``ellipticity condition\textquotesingle\textquotesingle\ is the \textbf{Legendre-Hadamard} condition
\begin{equation}\label{eq:legendre-hadamard-positivity-condition}
\eta^I\xi_kA_{kl}^{IJ}\eta^J\xi_l>0
\end{equation}
for every $\bm{x}\in\Omega$ and for every nonzero $\bm{\eta}\in\mathbb{R}^N$ and $\bm{\xi}\in\mathbb{R}^n$. The uniform version states that there exists $\theta>0$ such that for every $\bm{x}\in\Omega$
\begin{equation}\label{eq:legendre-hadamard-uniform-condition}
\eta^I\xi_kA_{kl}^{IJ}\eta^J\xi_l>\theta|\bm{\eta}|^2|\bm{\xi}|^2
\end{equation}
for every nonzero $\bm{\eta}\in\mathbb{R}^N$ and $\bm{\xi}\in\mathbb{R}^n$. These conditions turn out to be more physically reasonable than \eqref{eq:elliptic-system-full-matrix-positivity} or \eqref{eq:elliptic-system-full-matrix-uniform-bound} for many problems in elasticity. Note that \eqref{eq:legendre-hadamard-positivity-condition} and \eqref{eq:legendre-hadamard-uniform-condition} are much weaker than the corresponding conditions \eqref{eq:elliptic-system-full-matrix-positivity} and \eqref{eq:elliptic-system-full-matrix-uniform-bound}. (The inequalities have to hold only for rank-$1$ $N\times n$ matrices.) Despite this, \eqref{eq:legendre-hadamard-positivity-condition} and \eqref{eq:legendre-hadamard-uniform-condition} are sometimes referred to as \emph{strong ellipticity} conditions. As this example shows, the reader should be forewarned that the nomenclature surrounding elliptic systems does not necessarily make sense. More importantly, there is not universal agreement regarding these definitions. In reading the literature one needs to be careful to note the definitions various authors use.
